\section{Database}

\subsection{Theory}

\subsubsection{Introduction}

A database can be defined as a collection of related data items within a specific business.
DBMS are sotwares that define, create, use and maintain a database, if we combine a DBMS with a database we get a database system
The majority of DBMS softwares rely on SQL systems. 

\vspace{10pt}

To work with databases we need to model them to understand how they work. Modelling means representing reality in a 
scaled version. The same with databases we connect information to represent how relationships in the world exist.
Based on the business needings we design the database in a logical fashion and we translate it to 
a our language of choice. The step is passing from conceptual design to logical design. The first is DBMS independent while the second is DBMS dependent.
The process is:

\begin{enumerate}
    \item Requirement collection and analysis
    \item Conceptual design
    \item Logical design
    \item Physical design
\end{enumerate}

\vspace{10pt}

An SQL database is based on entities and relationships. Entities are business concepts
with an unambiguous meaning to a particular set of users. Entities have certain characteristics called attributes.
The rows are the entities (also called instances). The columns are the attributes. An association 
between one or more entities is called relationship, it defines how tables are related to each other. 

\vspace{10pt}

Another useful concept is the one of cardinality, how many instances of the entities can be associated with each other. 

We can have
 \begin{itemize}
    \item 1 to 1
    \item 1 to many
    \item Many to many
    \item None to many
    \item 1 to None
\end{itemize}

\vspace{10pt}

An important step in the design is the one related to data normalization, we want to avoid redundant data. We split 
the tables into smaller tables to avoid anomalies during insertion, deletion or updates. The 
procedure follows a step by step transformation. The database is more complex from the structural point of view but can be easier to manage.

We have 3 possible normal forms.

\begin{itemize}
    \item First NF \ra Each cell in the table can have only one value, never a list of values 
    \item Second NF \ra An entity type is in 2NF when it is in 1NF and when all of its non-key
attributes are functionally dependent on its primary key. Features should be fully dependent on the entire primary key
    \item Third NF \ra An entity type is in 3NF when it is in 2NF, and non-key attributes are
directly (non-transitively) dependent on the primary key.
In other words: non-prime attributes must be functionally dependent on the
key(s), but they must not depend on another non-prime attribute.

\end{itemize}

\vspace{10pt}


\subsubsection{DBSM organization and structure}


All DBMS follow a certain structure made by certain elements

\includegraphics[width=16cm]{dbms structure.png}


\begin{itemize}
    \item \textbf{Connection manager} \ra
    \item \textbf{Security manager}\ra
    \item \textbf{DDL compiler}\ra
    \item \textbf{Utilities}\ra
    \item \textbf{DML compiler}\ra
    \item \textbf{Query parser}\ra
    \item \textbf{Query rewriter}\ra
    \item \textbf{Query optimizer}\ra
    \item \textbf{Query executor}\ra
    \item \textbf{Transaction manager}\ra
    \item \textbf{Buffer manager}\ra
    \item \textbf{Lock manager}\ra
    \item \textbf{Recovery manager}\ra
\end{itemize}

A DBMS needs to interact with various parties, such as a database designer, a
database administrator, an application, or even an end-user (Ex. APIs).

\vspace{10pt}

To design our databases we can use an ERD (Entity Relationship Diagram) notation.

Some of them can be Chen, IDEF1X, Bachman, Crow's Foot (we use this), Min-Max or UML.

\includegraphics[width=8cm]{structure erd.png}

\includegraphics[width=8cm]{basic crows.png}

An entity is represented by a rectangle, with its
name on the top. The name is singular. The attribute(s) that uniquely distinguishes
an instance of the entity is the identifier.

Relationships have two indicators. The first one refers to the maximum number of times that an instance of one
entity can be associated with instances in the related entity. The second one describes the minimum number of times one instance can
be related to others. It can be zero or one, and optional or mandatory.

\includegraphics*[width=6cm]{multiplicities.png}
\includegraphics*[width=6cm]{options.png}


Relations can be
\begin{itemize}
    \item One to one
    \item One to many
    \item Many to many
\end{itemize}

Since we are working with tables we need to pass from schemas to tables, we use keys (foreign and primary)

